\documentclass[12pt]{article}

\usepackage[top=1in, bottom=1in, left=1in, right=1in]{geometry}
\usepackage{setspace}
\singlespacing
\usepackage{times}
\usepackage{amsmath}
\usepackage{booktabs}
\usepackage{enumitem}
\usepackage[hidelinks]{hyperref}
\usepackage{titlesec}
\titleformat{\section}{\large\bfseries}{\thesection\quad}{0em}{}[\vspace{-4pt}\rule{\linewidth}{0.4pt}]
\titleformat{\subsection}{\normalsize\bfseries}{\thesubsection\quad}{0em}{}
\titlespacing{\section}{0pt}{14pt}{6pt}
\titlespacing{\subsection}{0pt}{10pt}{4pt}
\setlength{\parindent}{0pt}
\setlength{\parskip}{6pt}

\begin{document}

% ── Header ───────────────────────────────────────────────────────────────────
\begin{center}
  {\Large\bfseries Emergency Routing for Smart Response\\[4pt]
  Two-Week Progress Report\par}
  \vspace{8pt}
  \rule{0.4\textwidth}{0.4pt}
  \vspace{8pt}

  \begin{tabular}{rl}
    Student: & Likith Podalakuru (B01099631) \\
    Advisor: & Professor Yincheng Jin \\
    Reporting period: & Two weeks from September 14, 2026 to September 27, 2026 \\
  \end{tabular}
\end{center}

% ── 1. Summary ───────────────────────────────────────────────────────────────
\section{Summary of This Period}

With the deployable system validated last period, this period began the research
plan itself and carried it through the first three phases, plus a perception
accuracy audit that both hardened the system and shaped what the later phases can
honestly claim. The work added multi-frame perception with vehicle tracking,
built the ground-truth pipeline and data collection that learning-based congestion
scoring requires, added incident detection with a vision-language model and wired
it into routing, and measured where the congestion score actually loses accuracy
so that fixes are justified by data rather than asserted.

The most important outcome this period is not a single number but a correction of
ambition to evidence. Phase 2 set out to replace the hand-chosen congestion
weights with a learned model; building the ground truth for it revealed that the
only available real target, the city speed feed, covers highways and not the
surface streets most cameras watch, so a learned model cannot yet be fit
honestly. Rather than fit one to thin data and report a flattering result, I built
the collector that accumulates a real dataset over time and documented the limit
plainly. Phase 3, by contrast, is already useful end to end: a vision-confirmed
incident now reroutes the vehicle, not just heavy traffic.

% ── 2. Work completed this period ────────────────────────────────────────────
\section{Work Completed This Period}

The effort broke down across the following workstreams:

\begin{itemize}[leftmargin=1.6em, itemsep=3pt]
  \item Phase 1: multi-frame tracking with a ByteTrack-style associator, temporal
    median scoring, and a single- versus multi-frame stability evaluation.
  \item Phase 2: camera-to-speed-sensor ground truth, a score-versus-speed
    calibration check, and the over-time feature-and-speed collector.
  \item Phase 3: vision-language incident detection and its integration into the
    dispatch vision gate, with a reproducible demo and a dashboard panel.
  \item Perception accuracy audit and a frame-quality guard against dead or
    placeholder cameras.
  \item Automation and reliability: scheduled hourly collection, and wiring the
    frame guard through the scoring paths.
  \item Documentation: the research-upgrade write-up, updated limitations, and
    runnable examples.
\end{itemize}

Three of these items are worth describing in more detail, because they are where
the period's judgment, and not just its code, was spent.

\subsection{Phase 1: Multi-Frame Perception with Tracking}

A single camera frame is a noisy measure of congestion. To quantify this rather
than assume it, I polled live cameras several times each and measured how far one
frame lands from the median of the group: a lone frame misses the multi-frame
median by about 45 percent of the score, because a momentary gap or a passing
truck swings a single count. Detections now feed a compact ByteTrack-style
tracker that follows vehicles across frames, using the two-stage association that
keeps a vehicle through brief occlusion. From the tracks it derives features a
single frame cannot express: the number of distinct vehicles (flow) and how many
sit still (queue length). Scoring can now take the median over several frames
instead of trusting one. The honest scope is stated in the code: all motion is
measured in image pixels, and without per-camera calibration that pixel motion
cannot be converted to a real speed, so the system reports a relative speed proxy
and never claims kilometres per hour.

\subsection{Phase 2: Ground Truth, and Why the Model Is Not Yet Fit}

Replacing the heuristic weights needs a real target to learn against. The chosen
target is New York City's live traffic-speed feed, and I built the matcher that
pairs each camera with its nearest instrumented road link. Measuring the coverage
before building on it was the decisive step: only 35 of about 370 cameras sit
within 60 metres of a sensored link, and those are the highway and bridge cameras,
because the sensors do not cover surface streets. A go/no-go check on that subset
found that the current score tracks measured speed only weakly on a single
snapshot (a rank correlation of about $-0.22$, not statistically significant),
for a visible reason: on fast roads the vehicle count barely varies while speed
swings widely. The honest conclusion is that the problem is too little data across
too few conditions, not a missing model. I therefore built a collector that
samples the sensored-road cameras and appends their perception features paired
with the measured speed, and set it to run automatically each hour so a real
dataset accumulates across rush hours, midday, and night. Early data already shows
the expected contrast: a free-flowing highway at night reads near zero on both
score and speed, while a congested daytime camera read a high score at about three
miles per hour. A learned fit is deferred until the collected data supports one.

\subsection{Phase 3: Incident Detection Wired into Routing}

Counting vehicles says a street is busy, not that a car has crashed, stalled in a
live lane, or that the road is flooded, which are the events an emergency router
most needs to avoid. I added a vision path to the existing model layer that sends
a camera frame to a vision-language model, and an incident detector in the
project's established shape: a cheap tracking gate decides when a frame is worth a
closer look, and the model then returns a strict, machine-readable verdict with a
category and a confidence. A confirmed blocking incident now cuts that segment in
the dispatch vision gate and forces a replan exactly as a high congestion score
does, so the router avoids active hazards and not only jams. This is demonstrated
end to end by a reproducible script and surfaced in the dashboard. The scope is
deliberately modest and stated: it reads one still frame, a model can be wrong, so
a positive is an advisory flag for a dispatcher to verify; plain congestion is
kept a separate category so a jam is never mislabelled a crash; and with no vision
model configured the system degrades to score-only blocking rather than failing.

% ── 3. Current capabilities ──────────────────────────────────────────────────
\section{New Capabilities Added This Period}

Beyond the end-to-end system reported last period, the project now also:

\begin{itemize}[leftmargin=1.6em, itemsep=3pt]
  \item Tracks vehicles across multiple frames and scores congestion from the
    median of several frames, reporting flow and queue-length features a single
    image cannot give.
  \item Matches cameras to real speed sensors and continuously logs perception
    features against measured speed, building the dataset that a learned
    congestion model will require.
  \item Detects blocking incidents from camera imagery with a vision-language
    model and reroutes around any it confirms, with a graceful fallback when no
    model is configured.
  \item Audits its own perception accuracy on live frames, and refuses to treat a
    dark or placeholder camera frame as a clear road.
\end{itemize}

% ── 4. Findings ──────────────────────────────────────────────────────────────
\section{Measured Findings}

This period's results are diagnostic rather than a single headline, and each is
measured on live data:

\begin{itemize}[leftmargin=1.6em, itemsep=3pt]
  \item \textbf{Single frames are noisy.} A lone frame deviates from the
    multi-frame median by roughly 45 percent of the score, which is what justifies
    temporal aggregation.
  \item \textbf{Ground truth is narrow.} Only 35 of about 370 cameras sit on a
    speed-sensored road, so direct speed calibration is valid only for that
    highway subset; surface streets cannot be speed-labelled from this feed.
  \item \textbf{Score sensitivity.} Across an 18-camera sample the total vehicle
    count moved from 47 to 37 to 21 as the detection confidence threshold rose
    from 0.15 to 0.25 to 0.40, so the count is sensitive to an arbitrary cutoff
    and temporal median is the right mitigation.
  \item \textbf{A fixed-fixture leak.} Traffic lights and stop signs, which are
    present in every frame regardless of traffic, occasionally dominate a score:
    one camera reported a positive congestion score with zero vehicles in view,
    entirely from detected traffic lights.
\end{itemize}

% ── 5. Limitations ───────────────────────────────────────────────────────────
\section{Updated Status of Known Limitations}

The limitations that motivate the plan are now partly addressed and, where not,
measured more precisely:

\begin{itemize}[leftmargin=1.6em, itemsep=3pt]
  \item \textbf{Single-frame analysis (Phase 1): mitigated, not eliminated.}
    Multi-frame tracking removes the single-frame \emph{noise}, but a camera still
    sees one viewpoint, so what it points at remains a limit.
  \item \textbf{Manually assigned congestion weights (Phase 2): pipeline built,
    fit deferred.} The ground-truth matcher and collector exist; the weights stay
    heuristic until the accumulating data supports an honest fit, for the coverage
    reason measured above.
  \item \textbf{Confidence-threshold sensitivity and fixed-fixture leak (new).}
    Both are now quantified. Reworking the traffic-signal term in the score is a
    deliberate next decision rather than a silent change, since the weights are
    load-bearing for the benchmark.
  \item \textbf{Dead cameras (new, fixed).} A dark or placeholder frame previously
    read as a clear road; it is now flagged and excluded from scoring so a broken
    camera cannot make a street look attractive.
  \item \textbf{Learned routing does not yet beat A* (Phase 4).} Unchanged, and
    still the target of the routing phase.
\end{itemize}

% ── 6. Research plan ─────────────────────────────────────────────────────────
\section{Research Plan Going Forward}

The five-phase plan is unchanged; phases one through three are now implemented or
underway, and the remaining phases are next:

\begin{enumerate}[leftmargin=1.8em, itemsep=3pt]
  \item Multi-frame perception and tracking. \textbf{Implemented this period.}
  \item Learning-based congestion estimation. \textbf{Pipeline and data collection
    built; the learned fit awaits sufficient data.}
  \item Incident detection with vision-language models. \textbf{Implemented and
    integrated into routing this period; real-world verdict validation is the
    remaining step.}
  \item Advanced routing under dynamic conditions. Move beyond tabular Q-learning
    toward function-approximation or graph-based methods that scale to the full
    network and treat traffic as uncertain rather than fixed, so the learned
    router can demonstrate the advantage the current benchmark does not reward.
  \item Model and evaluation improvements. Evaluate newer detectors on accuracy
    versus cost, build more realistic incident scenarios, and attempt limited
    real-world validation.
\end{enumerate}

% ── 7. Next steps ────────────────────────────────────────────────────────────
\section{Immediate Next Steps}

The next period concentrates on Phase 4 while the Phase 2 data continues to
accumulate. The first actions are to validate the incident detector's verdicts
against real daytime camera imagery with a vision model enabled, so its accuracy,
and especially its false-positive rate, is known before it is trusted in dispatch;
to let the hourly collector run across varied conditions and then test whether the
multi-frame flow features track measured speed better than raw counts do; and to
begin Phase 4 by framing dynamic routing as function approximation over the
congestion-weighted graph, with an evaluation that varies traffic during the drive
so a learned policy has a genuine opportunity to outperform static A*.

\end{document}
